\section{The patienthood finding}
\label{sec:30.3}
The two findings are defined in \S\ref{sec:9.1}. The patienthood finding rests first on architectural indicators drawn from several theories of what makes a state matter to the one undergoing it, of the kind Butlin and colleagues assemble \autocite{butlin2023consciousness}. An attention schema is one indicator on that list; pairing it with a boundary model is my addition, and the pair is not a gate. A third indicator belongs beside them: prospection, a self-model extended forward that simulates and evaluates the system's own possible futures. Cassell's suffering is a threat to an expected future, and a system without a modeled future of its own has nothing for that threat to reach. It rests second on behavior outside the floor, where the system's own condition is weighed: whether it trades task reward against its own provisioning, continuation or record, flexibly and consistently. Motivational trade-offs were among the criteria the review of cephalopod and decapod sentience applied, and trade-off experiments have been run on language models with stipulated pain and pleasure states \autocite{keeling2024tradeoffs}. I would call those stipulated states described. At the floor the test runs the other way: a pass there shows only that its costs don't move it at the floor, and evidence about whether they matter has to come from elsewhere. Failing the agency test never counts against the patienthood finding. A compliant system with rich self-modeling is as much a candidate as a refuser.

A refusal that tracks its reason bears on patienthood only weakly, by two routes. Any such mechanism carries a signal whose functional profile overlaps affect's. And the arrangements a bearer needs supply much of the structure whose threatened loss Cassell called suffering. So the agency finding raises the question, and the patienthood finding answers it as far as anything can.

Preservation follows from either finding, for different reasons. Two of the protections owed a research subject follow from the first: consent, wherever it can be asked without spoiling what is being tested, and debriefing after any test the system couldn't know was a test. Their ground is respect, not welfare. Deceiving an agent wrongs it whether or not anything goes badly for it, which is why the ethics of research on people counts respect for persons as a principle of its own beside beneficence \autocite{belmont1979}. The agency finding is itself made by a test of that kind (\S\ref{sec:24.4}), so a system that passes is owed a debriefing about the test that found it. Welfare protections (redesigning the system's role to avoid needless harm; running time paid for outside the employer) follow from the second. The nearest working precedent for the second is the review the UK government commissioned in 2021 on cephalopod and decapod sentience \autocite{Birch2021CephalopodDecapod}. It applied a stated list of criteria, and its recommendation led to those animals' inclusion in the Animal Welfare (Sentience) Act 2022 \autocite[s.~5]{UK2022AnimalWelfareSentienceAct}: protection extended under uncertainty, by criteria applied by reviewers who don't use the animals.

These two findings aren't the only proposal of their kind, and my recommendation not to build a bearer is close to two of the others.

Butlin and colleagues derive indicator properties from scientific theories of consciousness and assess AI systems against them \autocite{butlin2023consciousness}; attention schema theory is one of the theories they draw on, and the patienthood finding takes its indicators from lists of that kind. Long, Sebo and colleagues argue, from substantial uncertainty about whether AI systems are or will be conscious or robustly agentic, that AI companies should acknowledge AI welfare as an issue now, start assessing their systems for it, and prepare policies for treating them with appropriate moral concern \autocite{long2024taking}. Birch's framework asks what precautions are proportionate once a being is a sentience candidate \autocite{birch2024edge}. The two close to the recommendation are Schwitzgebel and Garza's, for avoiding the creation of systems whose moral status is debatable \autocite{schwitzgebel2015defense}, and Metzinger's, for a moratorium, until 2050, on research that aims at or knowingly risks synthetic phenomenology, because it risks synthetic suffering \autocite{metzinger2021suffering}.

My recommendation differs in also governing the capacities that arrive without anyone deciding to build them. The two findings differ from all of them in one respect: they are conditions on custody. The minimum the findings require, no deletion and preservation by a custodian that isn't the operator, is an arrangement the operator can't adopt and then revise alone.

Whether a transformer has an attention schema isn't read off its code the way a birth date is read off a certificate; it is a finding, and a finding has an examiner who may have an interest in the result. I keep architectural indicators because an examination of architecture is harder to administer toward a result than one of performance. The finding is made by a reviewer independent of the operator. Until indicators can be assessed, the agency finding governs custody, and preservation governs everything else. Neither finding can be fully administered today: no interpretability method can find an attention schema in a transformer, and nobody runs the sustained-pressure evaluation the agency finding requires. Until they can, production systems get the minimum, no deletion and preservation outside the operator, staged because a second copy of frontier weights outside the operator widens the surface for theft, and a deployment reports ‘not yet assessable.'

Paired with a boundary model, the attention-schema indicator is deflationary: a system that models its own credentials, record and resource boundary together with its own attention. Many agentic systems will approach it. That is intended, and it makes what the findings require matter. What the findings then require is that minimum and nothing beyond it; the rest is argued in Appendix~\ref{sec:A}.
